\chapter{Thiết bị nhúng ESP32}

\section{Phần cứng}

Thiết bị sử dụng board ESP32-CAM tích hợp camera OV5640 và microphone INMP441
kết nối qua giao thức I2S.

\begin{table}[htbp]
\centering
\caption{Thông số phần cứng ESP32}
\label{tab:esp32-hardware}
\begin{tabular}{>{\raggedright\arraybackslash}p{0.3\linewidth} >{\raggedright\arraybackslash}p{0.6\linewidth}}
\toprule
\rowcolor{headerblue} \textcolor{white}{\textbf{Thành phần}} & \textcolor{white}{\textbf{Thông số}} \\
\midrule
\rowcolor{rowgray} Board & ESP32-CAM với PSRAM \\
Camera & OV5640, JPEG, QVGA (320$\times$240) mặc định, hỗ trợ đến UXGA (1600$\times$1200) \\
\rowcolor{rowgray} Microphone & INMP441 (I2S digital), 16 kHz, mono \\
Kết nối & WiFi 802.11 b/g/n \\
\rowcolor{rowgray} Giao thức & TCP (camera port 3001, audio port 3002) \\
\bottomrule
\end{tabular}
\end{table}

\subsection{Cấu hình chân kết nối}

\textbf{Camera OV5640} sử dụng bus dữ liệu 8-bit song song:
D0--D7 (GPIO 5, 18, 19, 21, 36, 39, 34, 35), XCLK (GPIO 0), PCLK (GPIO 22),
VSYNC (GPIO 25), HREF (GPIO 23), SDA (GPIO 26), SCL (GPIO 27), PWDN (GPIO 32).

\textbf{Microphone INMP441} kết nối qua I2S:
WS -- Word Select (GPIO 2), SCK -- Serial Clock (GPIO 14), SD -- Serial Data (GPIO 15),
sử dụng I2S port 1.

\section{Kiến trúc Dual-Core}

ESP32 có 2 nhân xử lý, được phân công như sau:

\begin{table}[htbp]
\centering
\caption{Phân công xử lý dual-core trên ESP32}
\label{tab:dual-core}
\begin{tabular}{>{\centering\arraybackslash}p{0.12\linewidth} >{\raggedright\arraybackslash}p{0.25\linewidth} >{\raggedright\arraybackslash}p{0.5\linewidth}}
\toprule
\rowcolor{headerblue} \textcolor{white}{\textbf{Core}} & \textcolor{white}{\textbf{Nhiệm vụ}} & \textcolor{white}{\textbf{Chi tiết}} \\
\midrule
\rowcolor{rowgray} Core 0 & Audio Capture Task & Đọc I2S liên tục từ INMP441, chuyển đổi 32-bit→16-bit PCM, stream qua TCP port 3002. Stack: 16 KB. \\
Core 1 & Camera + Main Loop & Chụp JPEG từ OV5640, gửi qua TCP port 3001, xử lý lệnh điều khiển camera từ server. \\
\bottomrule
\end{tabular}
\end{table}

Việc tách riêng hai core đảm bảo luồng audio chạy liên tục không bị gián đoạn
bởi camera capture, tránh hiện tượng buffer underrun.

\section{Pipeline thu âm thanh}

\subsection{Cấu hình I2S}

\begin{table}[htbp]
\centering
\caption{Cấu hình I2S cho microphone INMP441}
\label{tab:i2s-config}
\begin{tabular}{>{\raggedright\arraybackslash}p{0.35\linewidth} >{\raggedright\arraybackslash}p{0.55\linewidth}}
\toprule
\rowcolor{headerblue} \textcolor{white}{\textbf{Tham số}} & \textcolor{white}{\textbf{Giá trị}} \\
\midrule
\rowcolor{rowgray} Chế độ & I2S\_MODE\_MASTER | I2S\_MODE\_RX \\
Tần số lấy mẫu & 16.000 Hz \\
\rowcolor{rowgray} Bit depth đầu vào & 32-bit (I2S native) \\
Bit depth đầu ra & 16-bit PCM (sau chuyển đổi) \\
\rowcolor{rowgray} Kênh & Mono (chỉ kênh trái) \\
Số DMA buffer & 16 (tăng từ 8 để ổn định) \\
\rowcolor{rowgray} Kích thước buffer & 1024 mẫu (64 ms mỗi buffer) \\
Throughput & 32 KB/s ($\approx$ 256 kbps) \\
\bottomrule
\end{tabular}
\end{table}

\subsection{Chuẩn hóa tín hiệu}

Microphone INMP441 có độ phân giải 18-bit, đóng gói trong word 32-bit với dữ liệu
nằm ở các bit cao. Quá trình chuẩn hóa:

\begin{enumerate}
\item Đọc mẫu 32-bit raw từ I2S DMA buffer.
\item Dịch phải 14 bit: \texttt{val = raw\_samples[i] >> 14}.
\item Cắt ngưỡng vào phạm vi 16-bit signed: $[-32768, 32767]$.
\item Lưu dưới dạng \texttt{int16\_t} để truyền qua TCP.
\end{enumerate}

\section{Pipeline streaming video}

Camera OV5640 được cấu hình chụp JPEG với double buffering (2 frame buffer).
XCLK mặc định 12 MHz, chất lượng JPEG có thể điều chỉnh (10--63).

Mỗi frame được truyền qua TCP port 3001 với giao thức:
\texttt{[Length (4 bytes LE)][JPEG Data (variable)]}.
Buffer tĩnh 250 KB được cấp phát một lần trên PSRAM (hoặc heap) để tránh
phân mảnh bộ nhớ.

\section{Giao thức điều khiển camera}

Trình duyệt gửi lệnh điều khiển qua WebSocket đến server trung gian,
server chuyển tiếp qua TCP đến ESP32 với format:
\texttt{[0xA5][ID (1 byte)][Value (1 byte)]}.

\begin{table}[htbp]
\centering
\caption{Các lệnh điều khiển camera chính}
\label{tab:camera-commands}
\begin{tabular}{>{\centering\arraybackslash}p{0.08\linewidth} >{\raggedright\arraybackslash}p{0.3\linewidth} >{\raggedright\arraybackslash}p{0.45\linewidth}}
\toprule
\rowcolor{headerblue} \textcolor{white}{\textbf{ID}} & \textcolor{white}{\textbf{Tham số}} & \textcolor{white}{\textbf{Phạm vi}} \\
\midrule
\rowcolor{rowgray} 1 & JPEG Quality & 10--63 (thấp = chất lượng cao hơn) \\
2 & Frame Size & QQVGA đến UXGA (14 mức) \\
\rowcolor{rowgray} 3--5 & Brightness, Contrast, Saturation & $-2$ đến $+2$ \\
6--9 & AWB, AEC, AGC, AEC2 & 0/1 (bật/tắt) \\
\rowcolor{rowgray} 16--17 & H-Mirror, V-Flip & 0/1 \\
20 & XCLK Frequency & 6/8/12/16/24 MHz \\
\rowcolor{rowgray} 21 & Frame Buffer Count & 1--4 \\
\bottomrule
\end{tabular}
\end{table}

ESP32 phản hồi cài đặt hiện tại qua packet với magic header \texttt{0xBEEF}
theo format: \texttt{[0xBE][0xEF][ID][Value]...} cho tất cả 21 tham số.

\section{Server trung gian (Node.js)}

Server trung gian chạy trên Bun/Express đóng vai trò cầu nối giữa ESP32 và
trình duyệt:

\begin{itemize}
\item \textbf{TCP Camera Server (port 3001):} Nhận JPEG frame từ ESP32,
  parse 4-byte length header, broadcast qua WebSocket \texttt{/camera} đến browser.
\item \textbf{TCP Audio Server (port 3002):} Nhận stream 16-bit PCM từ ESP32,
  duy trì byte-alignment, broadcast qua WebSocket \texttt{/audio} đến browser.
\item \textbf{HTTP API Proxy:} Chuyển tiếp request từ browser đến các AI services:
  \begin{itemize}
  \item \texttt{POST /api/process-voice} → Denoise Service → STT Service
  \item \texttt{POST /api/gen-description} → Voice Rephraze Service
  \item \texttt{POST /api/classify} → Classification Service
  \end{itemize}
\end{itemize}

Giao diện web hiển thị video streaming real-time, biểu đồ amplitude âm thanh
(30 bars), spectrogram trước/sau khử nhiễu (WaveSurfer.js), kết quả nhận dạng
và phân loại sản phẩm.
